\begin{table}[htbp]
  \centering
  \caption{Os comandos de telemetria do kw. Cada um chama o kw\_tele
    \citep{kwfork_telemetry}.}
  \label{tab:comandos-telemetria}
  \sbox0{\texttt{kw telemetry <comando>}}%
  \providecommand{\kwopcao}{}\renewcommand{\kwopcao}[1]{\hspace{1.5em}\texttt{-{}-#1}}%
  \begin{tabular}{@{}l>{\raggedright\arraybackslash\hyphenpenalty=10000}p{\dimexpr\textwidth-\wd0-2\tabcolsep\relax}@{}}
    \toprule
    \texttt{kw telemetry <comando>} & Descrição \\
    \midrule
    \texttt{upload}             & Envia agora os eventos pendentes \\
    \kwopcao{dry-run}           & Mostra o JSON do próximo lote, sem enviar \\
    \cmidrule(l{1.5em}r{1.5em}){1-2}
    \texttt{consent}            & Mostra o estado; se desligada e há terminal, pergunta \\
    \cmidrule(l{1.5em}r{1.5em}){1-2}
    \texttt{opt-in}             & Liga a telemetria, sem os eventos já registrados \\
    \kwopcao{include-history}   & Liga a telemetria, com os eventos já registrados \\
    \cmidrule(l{1.5em}r{1.5em}){1-2}
    \texttt{opt-out}            & Desliga a telemetria \\
    \cmidrule(l{1.5em}r{1.5em}){1-2}
    \texttt{info}               & Mostra o que é coletado e o estado dos envios \\
    \kwopcao{export}            & Exporta em JSON o que já foi enviado \\
    \cmidrule(l{1.5em}r{1.5em}){1-2}
    \texttt{config}             & Mostra ou muda os tempos de envio e o tamanho do lote \\
    \bottomrule
  \end{tabular}
\end{table}
